\subsection{The Diamond-Lift Closure Conjecture}

First stated in \texttt{dev/content/notes/2026/07/02/} at \texttt{metroplat.emsenn.net}
(``Diamond-Lift Closure Conjecture''). Where \cref{thm:hole-diamond-kinematic-space} packages the
causal diamonds of a causal set $C$ as their own sequencing causal set
$K(C) := \texttt{HoleDiamondKinematicSpace}(C)$, this conjecture asks whether a FOUR-coloured
post-hole tile over $C$ (\texttt{PostHoleTile}, \cref{ch:four}'s recurring coherence cell realised
as a coloured tile) lifts, along the diamond's upper endpoint, to a FOUR-coloured post-hole tile
over $K(C)$.

Building a real instance first surfaced two frontier gaps in the pre-existing \texttt{PostHoleTile}
carrier itself, neither previously blueprinted: \texttt{HoleTileAdjacency} was reflexive
(\texttt{CausalAdjacent} always holds at the diagonal), making every nonempty-hole colouring
obligation unsatisfiable outright; and it quantified over the whole ambient type $E$ rather than
the causal set's finitely many listed events, so over an unbounded $E$ (e.g. \texttt{Articulating})
the computed hole region is cofinite and its comparability graph contains chains longer than four,
again unsatisfiable. Both are now closed.

\begin{theorem}[Post-hole tile adjacency is a genuine finite graph]\label{thm:post-hole-tile-adjacency-fixed}\lean{Obsmetrologia.CausalSets.BlackHoles.HoleTileAdjacency, Obsmetrologia.CausalSets.BlackHoles.holeTileAdjacency_ne, Obsmetrologia.CausalSets.BlackHoles.holeTileAdjacency_mem_left, Obsmetrologia.CausalSets.BlackHoles.holeTileAdjacency_mem_right}\leanok
\texttt{HoleTileAdjacency} now requires \texttt{x ≠ y} (excluding the unavoidable self-loop
\texttt{CausalAdjacent} carries at every point) and \texttt{x, y ∈ C.events} (excluding
comparability with events the causal set never actually lists). Both guards are load-bearing:
without the first, \texttt{HoleTileColoring} demands \texttt{color x ≠ color x} at every hole
event; without the second, the comparability graph on an unbounded causal set (any chain longer
than four, e.g. an articulation chain) is not four-colourable at all. Neither guard was previously
present, and \texttt{PostHoleTile} had no instance anywhere in the library before this pass.
\end{theorem}

\begin{theorem}[A genuine post-hole tile witness]\label{thm:post-hole-tile-witness}\lean{Obsmetrologia.CausalSets.BlackHoles.dlCausalSet, Obsmetrologia.CausalSets.BlackHoles.dlReadout, Obsmetrologia.CausalSets.BlackHoles.dlColor, Obsmetrologia.CausalSets.BlackHoles.dlColoring, Obsmetrologia.CausalSets.BlackHoles.dlAnchorCore, Obsmetrologia.CausalSets.BlackHoles.dlPostHoleTile}\leanok\uses{thm:post-hole-tile-adjacency-fixed}
\texttt{DLChain} is a genuinely Mark-downstream carrier: three real \texttt{Articulating} values
(one \texttt{markStepCausalSet} top, applied twice over a bare uttered floor), packaged as a
finite subtype rather than left over the unbounded \texttt{Articulating} type, so
\texttt{Hole.lean}'s future-core machinery -- whose closure/connectivity quantifiers range over
the whole ambient type -- applies to what it was built for. With the readout at the floor, the
computed hole is exactly $\{\texttt{level1}, \texttt{top}\}$; \texttt{dlColor} properly colours
its sole edge, and \texttt{dlAnchorCore} anchors the whole hole at \texttt{top}. \texttt{dlPostHoleTile}
is the resulting witness -- the first inhabited \texttt{PostHoleTile} anywhere in the library.
\end{theorem}

\begin{theorem}[The readout exterior is connected too]\label{thm:readout-exterior-connected-witness}\lean{Obsmetrologia.Alexandrov.isConnected_singleton, Obsmetrologia.CausalSets.BlackHoles.dlReadout_asProp_eq_singleton, Obsmetrologia.CausalSets.BlackHoles.dlReadoutConnected, Obsmetrologia.CausalSets.BlackHoles.dlReadoutExteriorConnected}\leanok\uses{thm:post-hole-tile-witness}
\texttt{ReadoutExteriorConnected} is \texttt{HoleRegionConnected}'s sibling in \texttt{Hole.lean}
and had the same defect: stated, reasoned about generically, never once witnessed. It needed no new
combinatorics, only a missing base case in the Alexandrov-connectedness toolkit itself --
\texttt{isConnected\_upperSet} was the sole generator of concrete connected sets in the whole
library, and a singleton readout like \texttt{dlReadout} (\cref{thm:post-hole-tile-witness}) is not
an up-set. \texttt{isConnected\_singleton} adds that base case upstream, in
\texttt{Cohomology.AlexandrovHusainov} alongside \texttt{isConnected\_upperSet} itself: a one-point
set admits no nontrivial separation, so it is connected outright. \texttt{dlReadoutConnected}
transports it along \texttt{dlReadout}'s Boolean-to-Prop reading, and
\texttt{dlReadoutExteriorConnected} then applies the existing
\texttt{readoutExterior\_connected\_of\_connected\_listed\_readout} directly -- the missing piece
was upstream, not downstream.
\end{theorem}

\begin{theorem}[The diamond upper-endpoint map is monotone, and hole membership pulls back]\label{thm:diamond-lift-general}\lean{Obsmetrologia.CausalSets.BlackHoles.upperPoint_monotone, Obsmetrologia.CausalSets.BlackHoles.holeRegion_lift_pullback}\leanok\uses{thm:hole-diamond-kinematic-space}
For any causal set $C$, the map $p \mapsto p.\text{upper}.\text{point}$ from $K(C)$ to $C$ is
monotone: a sequencing edge (or equality) between diamond points gives a causal inequality between
their upper endpoints. Consequently, pulling a readout back along this map
($R^\diamond(p) := R(p.\text{upper}.\text{point})$) pulls hole membership back too: if the upper
endpoint of a diamond point lies in the base hole region, the diamond point itself lies in the
lifted hole region on $K(C)$. Both hold for an arbitrary causal set, not just the concrete witness
below.
\end{theorem}

\begin{theorem}[The lifted colouring is proper]\label{thm:diamond-lift-coloring}\lean{Obsmetrologia.CausalSets.BlackHoles.dlReadoutD, Obsmetrologia.CausalSets.BlackHoles.dlColorD, Obsmetrologia.CausalSets.BlackHoles.dlHoleRegionD_true, Obsmetrologia.CausalSets.BlackHoles.dlColoringD}\leanok\uses{thm:diamond-lift-general, thm:post-hole-tile-witness}
Lifting \texttt{dlReadout}/\texttt{dlColor} along the upper endpoint puts every diamond point of
the witness chain into the lifted hole region (\cref{thm:diamond-lift-general} applied to
\texttt{dlDiamondPoint\_upper\_ne\_dlFloor}), and \texttt{dlColorD} properly colours the lifted
adjacency graph: every sequencing step over this three-event chain runs from the diamond ending at
\texttt{level1} to the diamond ending at \texttt{top} -- the chain's only direct edge -- so the
base colouring's one distinguishing fact carries over exactly.
\end{theorem}

\begin{theorem}[The lift's non-closure is a provable collapse, not a gap in the proof]\label{thm:diamond-lift-collapse}\lean{Obsmetrologia.CausalSets.BlackHoles.dlNoSuccessor_of_upper_dlTop, Obsmetrologia.CausalSets.BlackHoles.dlDPJump_ne_dlDPStep', Obsmetrologia.CausalSets.BlackHoles.dlDPJump_dlDPStep'_upperPoint_collapse, Obsmetrologia.CausalSets.BlackHoles.dlDPJump_dlDPStep'_colorD_collapse, Obsmetrologia.CausalSets.BlackHoles.dlDPJump_dlDPStep'_unrelated}\leanok\uses{thm:diamond-lift-coloring}
The obstruction to a connected future core is not merely observed in the witness; it is derived.
The upper-endpoint map $p \mapsto p.\text{upper}.\text{point}$ is not injective: the long-jump
diamond $\texttt{dlDPJump} = (\texttt{floor},\texttt{top})$ and the second-step diamond
$\texttt{dlDPStep'} = (\texttt{level1},\texttt{top})$ are two \emph{distinct} diamond points
(\texttt{dlDPJump\_ne\_dlDPStep'}, since they disagree at the lower endpoint) that collapse onto
the \emph{same} base event and hence the same lifted colour
(\texttt{dlDPJump\_dlDPStep'\_upperPoint\_collapse}, \texttt{dlDPJump\_dlDPStep'\_colorD\_collapse}).
This is not a harmless colour coincidence: \texttt{dlNoSuccessor\_of\_upper\_dlTop} proves that no
diamond point ending at \texttt{top} has \emph{any} sequencing successor at all (nothing is reached
from the top of this chain), and \texttt{dlDPJump\_dlDPStep'\_unrelated} uses it to show
$\texttt{dlDPJump}$ and $\texttt{dlDPStep'}$ are causally unrelated in \emph{both} directions in
$K(\texttt{dlCausalSet})$. So the long-jump diamond is a genuine isolated point of the lifted hole
region -- proved isolated, not merely absent from the one edge the base tile's hole carries.
\end{theorem}

\begin{theorem}[Flipping the endpoint closes the lift]\label{thm:diamond-lift-lower-closes}\lean{Obsmetrologia.CausalSets.BlackHoles.lowerPoint_monotone, Obsmetrologia.CausalSets.BlackHoles.holeRegion_lift_pullback_lower, Obsmetrologia.CausalSets.BlackHoles.DiamondPoint.cases_dl, Obsmetrologia.CausalSets.BlackHoles.dlReadoutDLower, Obsmetrologia.CausalSets.BlackHoles.dlColorDLower, Obsmetrologia.CausalSets.BlackHoles.dlHoleRegionDLower_iff, Obsmetrologia.CausalSets.BlackHoles.dlColoringDLower, Obsmetrologia.CausalSets.BlackHoles.dlFutureCoreDLower, Obsmetrologia.CausalSets.BlackHoles.dlPostHoleTileDLower}\leanok\uses{thm:diamond-lift-collapse}
The conjecture's phrasing (``an endpoint-reflecting lift'') names no preferred endpoint, and reading
$R^\diamond$/$\chi^\diamond$ off the diamond's \emph{lower} endpoint is dual and equally honest:
\texttt{lowerPoint\_monotone}/\texttt{holeRegion\_lift\_pullback\_lower} are the exact mirrors of
\cref{thm:diamond-lift-general}, general over any causal set (`$p.\text{lower} \le p.\text{upper}
\le q.\text{lower}$' whenever $p \prec^\diamond q$ supplies the missing step). On the witness chain
this flip changes everything: \texttt{DiamondPoint.cases\_dl} enumerates $K(\texttt{dlCausalSet})$'s
three diamond points exhaustively, and \texttt{dlHoleRegionDLower\_iff} shows the lower-pullback hole
collapses to the \emph{singleton} $\{\texttt{dlDPStep'}\}$ --- \texttt{dlDPStep} and
\texttt{dlDPJump} both have lower endpoint \texttt{floor}, so both reflexively reach the (lower-read)
readout and drop out. A singleton hole is a \texttt{PostHoleTile} for free: no adjacency pairs to
colour (\texttt{dlColoringDLower}, vacuous) and trivially connected
(\texttt{dlFutureCoreDLower}, via \texttt{isConnected\_singleton}). \texttt{dlPostHoleTileDLower}
is the resulting witness --- a \emph{fully closed} lift of \texttt{dlPostHoleTile}, on the same base
tile \cref{thm:diamond-lift-collapse} shows the upper-endpoint lift cannot close. The upper-endpoint
obstruction was real, but local to that specific lift, not to the diamond space itself.
\end{theorem}

\begin{theorem}[Diamond-Lift Closure Conjecture]\label{thm:diamond-lift-closure-conjecture}\notready\uses{thm:diamond-lift-collapse, thm:diamond-lift-lower-closes}
\emph{(Open frontier --- stated, not formalized.)} For every FOUR-coloured post-hole tile $T$ over
a causal set $C$, if the readout and colouring admit an endpoint-reflecting lift to $K(C)$ and the
lifted hole adjacency is the sequencing-controlled adjacency on diamond points, then the lift is
again a FOUR-coloured post-hole tile: $(C, R, \chi) \rightsquigarrow (K(C), R^\diamond, \chi^\diamond)$
is a closure operation on post-hole FOUR tiles. The colouring half is proved concretely for the
upper-endpoint lift (\cref{thm:diamond-lift-coloring}); \cref{thm:diamond-lift-collapse} proves that
lift's future core cannot close, by a genuine non-injective collapse onto an isolated point; and
\cref{thm:diamond-lift-lower-closes} proves the \emph{lower}-endpoint lift of the same base tile
closes completely. So the witness resolves to: existence of \emph{some} closing endpoint-reflecting
lift, not universal closure of \emph{every} one. The still-open content is exactly that existential
--- a general theorem (over an arbitrary causal set and post-hole tile, not just this three-event
witness) that at least one of the lower- or upper-endpoint lifts always closes, or a precise
structural criterion for which one does.
\end{theorem}
